\chapter{一套实用的 PCSEL 仿真工作流}
\label{ch:workflow}

\section*{学习目标}
\begin{itemize}
  \item 把前面二十余章建立的概念、方程和近似组织成一条真正可执行的研究工作流。
  \item 明确不同阶段该计算什么、向下一阶段交付什么，以及哪些结论在该阶段绝不能提前宣称。
  \item 在求解器无关的原则下，理解 Lumerical 与 COMSOL 分别更适合承担哪些计算，并掌握最小可执行建模过程。
\end{itemize}

\section*{先修知识}
建议已掌握 \cref{ch:method-compare,ch:epitaxial-optics,ch:electrical,ch:eto-coupling,ch:robustness,ch:dynamics,ch:metrics} 中关于方法分工、外延层、电学、热学、鲁棒性分析、动态稳定性和性能报告的内容。

\section*{本章路线图}
到这里，全书已经把 PCSEL 的核心物理链条铺开了。剩下的问题不再是“某个公式怎么推”，而是“研究时应该先做什么、后做什么、哪些结果能信、哪些结果还只是候选”。本章的目标就是把这条链变成工作流。我们先建立分层结论和数据账本，再逐阶段说明纵向模检查、周期单胞筛选、有限器件验证、电热建模、热-电-光回写以及鲁棒性分析的输入输出，随后专门讨论两条常见实现路径：以 Lumerical 为主的光学优先链，以及以 COMSOL 为主的多物理场优先链。最后给出停机判据和结果交付模板。

\section{本章如何与下一章例题和失败模式章配合使用}
如果只把本章当成“流程图章节”，读者很容易知道名词，却不知道每一步真正长什么样；如果只看下一章例题，又可能误以为那只是某一个特定结构的演示，而不是一般工作流的落地版本。更稳妥的读法是把 \cref{ch:workflow,ch:worked-example,ch:failure-modes} 连成一套使用：本章给出正向流程，\cref{ch:worked-example} 把它落实成一次完整作答，而 \cref{ch:failure-modes} 则告诉你每一步最常见的失误会长成什么样。

\begin{table}[htbp]
  \centering
  \footnotesize
  \caption{工作流、例题与失败模式的对应关系。}
  \label{tab:workflow-example-failure}
  \setlength{\tabcolsep}{4pt}
  \begin{tabularx}{\textwidth}{@{}p{1.3cm}p{3.0cm}Y Y@{}}
    \toprule
    阶段 & 本章任务 & 在 \cref{ch:worked-example} 中的落地位置 & 若这一步做错，在 \cref{ch:failure-modes} 中最接近的失败类型 \\
    \midrule
    A & 纵向模与层栈检查 & 例题 B 的“第一步：纵向模检查” & 把错误层栈背景带入后续所有计算；误用有效折射率近似 \\
    B & 周期单胞筛选 & 例题 B 的“第二步：单胞候选模筛选” & 把单胞开放结果直接当器件结论；误读被动排序 \\
    C & 有限器件全波验证 & 例题 B 的“第三步：有限阵列检查” & 有限尺寸建模不足；边界条件与几何定义不一致 \\
    D & 电注入与热建模 & 例题 B 的“第四步：简化电学/热学回写” & 把均匀增益当真实注入；忽略热回写 \\
    E/F & 工作窗口与鲁棒性 & 例题 B 的“中间结果汇总与结论分层” & 名义值最优但无工艺窗口；对工作窗口层级越权宣称 \\
    G & 实验闭环 & 例题 B 的“哪些问题必须交给实验闭环” & 用已拟合量冒充预测能力；缺少回验量 \\
    \bottomrule
  \end{tabularx}
\end{table}

因此，最推荐的顺序不是“先看完本章再说”，而是：先读本章建立流程骨架，再在 \cref{ch:worked-example} 中顺着这张表把每一步真正做一遍，最后用 \cref{ch:failure-modes} 把整条链反过来检查一遍。

\begin{table}[htbp]
  \centering
  \footnotesize
  \caption{工作流各阶段的典型计算资源量级。这里给的是项目排程意义下的数量级估计，而不是普适硬门槛。}
  \label{tab:workflow_resource_budget}
  \begin{tabularx}{\textwidth}{@{}p{1.1cm}p{3.3cm}p{2.6cm}Y@{}}
    \toprule
    阶段 & 主要任务 & 典型耗时量级 & 资源判断 \\ \midrule
    A & 纵向模与层栈检查 & 秒到分钟 & 二维模态问题，内存压力通常较小，适合大量参数初筛 \\
    B & PWEM / RCWA 单胞筛选 & 秒到分钟；做参数扫时可到小时 & 周期单胞便宜，但参数扫会迅速放大总成本 \\
    C & 有限阵列 FDTD / FEM & 小时到天 & 三维全波往往是第一道真正的算力门槛，常需多核、大内存或 GPU \\
    D & 电注入与热建模 & 分钟到小时 & 比全波便宜，但网格与非线性收敛会明显受几何复杂度影响 \\
    E/F & 热回写与鲁棒性分析 & 小时到天 & 单次求解未必最慢，但多轮迭代和多样本统计常使总成本最高 \\
    G & 实验闭环与校准 & 组织成本最高 & 计算时间未必最长，但参数整理、版本控制和回验最消耗研究日程 \\
    \bottomrule
  \end{tabularx}
\end{table}

\section{工作流的第一原则：先定义结论层级，再跑仿真}
仿真失败常常不是因为不会点按钮，而是因为在第一步就没有定义清楚自己打算得到哪一层结论。对 PCSEL 来说，最安全的做法是把结论分为四层：
\begin{enumerate}
  \item L1 结构候选层：只谈候选模、对称性、被动辐射趋势和纵向重叠；
  \item L2 光学阈值层：只谈开放结构中的阈值增益、偏振趋势、上下出光比例和有限尺寸模式排序；
  \item L3 电学映射层：开始谈注入分布、J-V、串联电阻、阈值电流区间的来源；
  \item L4 热-电-光工作层：开始谈工作窗口、稳定性、亮度、偏振稳定性与热漂移后的器件表现。
\end{enumerate}
如果还停留在 L1 或 L2，就绝不能提前宣称阈值电流、线宽窗口或高温稳定工作。这个分层不是文风问题，而是为了防止求解器输出和物理结论错位。

\begin{RigorousBox}
PCSEL 工作流中最重要的纪律是：\textbf{低层结论可以进入高层模型，高层结论不能倒过来由低层结果“脑补”出来}。被动 $Q$ 可以进入阈值增益估计，但被动 $Q$ 本身不能直接变成阈值电流或工作线宽。
\end{RigorousBox}

\section{工作流的第二原则：先建数据账本，再建几何}
在任何求解器启动之前，先建立参数账本 $\mathcal{D}$。它至少应记录四类信息：几何参数、材料参数、外延层参数和电热参数。每个参数都至少带四个元信息：单位、来源、适用温度或频段、当前不确定度。没有元信息的参数，不应直接进入主结果模型。

这个步骤看起来“管理味很重”，其实是研究效率的核心。因为 PCSEL 的仿真链跨越多个求解器，若不在一开始统一单位、坐标和参数来源，后面最容易发生的不是数值误差，而是参数错配。例如，光学模型用的是冷态折射率，电热模型却默认了工作温度；或者 RCWA 的单胞孔径与 CHARGE 的电极开窗尺寸根本不是同一个版本。一个看似微小的版本错配，足以让整条工作流失去可复核性。

\paragraph{把“数据账本”落实成真正可执行的 manifest}
为了让这条原则不只停在理念层，建议把参数账本直接写成一个可版本控制的 manifest。最小字段至少应包括：参数唯一 ID、显示符号、文件字段名、单位、数值、来源、适用波长/温度、被哪些阶段使用，以及当前不确定度。一个最小 YAML 模板可以长成下面这样：

\begin{verbatim}
parameter:
  id: MAT-GAAS-N-001
  name: n_GaAs
  symbol: n
  file_key: n_gaas
  unit: 1
  value: 3.42
  source: literature
  valid_wavelength_nm: [930, 990]
  valid_temperature_C: [20, 80]
  uncertainty: 0.02
  used_in:
    - stage_A_waveguide
    - stage_B_unitcell
    - stage_C_finite_array
  version: v1.2
  last_updated: 2026-03-01
\end{verbatim}

真正项目化时，这个 manifest 还应再加两列：\texttt{propagates\_to}（它会影响哪些输出量）和 \texttt{status}（direct-measured / literature / fitted / inferred）。这样一来，读者第一次独立搭建建模链时，就不会把“文献里抄来的冷态参数”和“实验反演出来的工作点参数”混写进同一个无版本表格。

\section{最小命令行复现骨架：即使使用商业软件，也应把文件流做成终端原生}
本书不把 GUI 按钮教程当主线，但这不等于 workflow 只能停在“软件名词表”。更稳妥的做法，是先把目录、参数文件、阶段脚本和输出制品固定成一条 terminal-native 骨架，再决定每一步内部到底由 Lumerical、COMSOL 还是其他求解器来完成。一个最小项目骨架可以写成：

\begin{verbatim}
pcsel_project/
  params.yaml
  01_layers/
  02_unitcell/
  03_finite_array/
  04_electrical_thermal/
  05_postprocess/
  run_stage_A.py
  run_stage_B.py
  run_stage_C.py
  run_stage_D.py
  run_stage_E.py
\end{verbatim}

\noindent 其中每一步至少应有固定的输入、输出、sanity checks 与失败回扫项：
\begin{itemize}
  \item Stage A：输入层栈与材料表，输出 $n_{\mathrm{eff}}$、$\Gamma_{\mathrm{QW}}$、$\Gamma_{\mathrm{loss}}$；sanity check 是模式主峰位置和量纲自洽。
  \item Stage B：输入单胞几何与层栈背景，输出候选模频段、被动损耗排序、上下辐射比例；sanity check 是 PWEM / RCWA 频段与对称性标签一致。
  \item Stage C：输入有限阵列几何，输出近场、远场、主模排序和尺寸收敛；sanity check 是尺寸/PML/网格继续加密后主结论不翻转。
  \item Stage D：输入外延、电极与接触参数，输出 $J(\rvec)$、$N(\rvec)$、$T(\rvec)$；sanity check 是边界条件、单位和功率守恒一致。
  \item Stage E：输入前四阶段结果，输出回写后的净模裕量、工作窗口和鲁棒性报告；sanity check 是结论层级不越权。
\end{itemize}

这套骨架的目的不是承诺“整本书都提供开源全复现脚本”，而是强制读者把文件流、参数流和结论层级固定下来。只要这一步做对，后续无论内部调用商业软件还是自写脚本，研究链条都会清楚得多。

\section{阶段 A：先做纵向模和层栈检查，而不是急着扫平面晶格}
PCSEL 的第一步不应该是直接扫单胞孔径，而应该是先确认外延层本身给出的纵向舞台是否合理。也就是说，先解层栈中的纵向模、有效折射率、量子阱重叠和损耗重叠。若这一关都没有过，后续平面内晶格优化大概率是建立在错误舞台上的。

在 Lumerical 路径里，这一阶段通常最适合用 \software{MODE/FDE} 之类的波导模求解器：建立二维横截面层栈，定义材料色散和损耗，求出目标偏振附近的导模，输出 $n_{\mathrm{eff}}$、场分布、$\Gamma_{\mathrm{QW}}$ 与与损耗层重叠积分。最小建模过程可以写成：
\begin{enumerate}
  \item 画出外延层截面并赋予折射率、掺杂相关吸收代理；
  \item 在工作波长附近做本征模求解；
  \item 提取目标模的 $n_{\mathrm{eff}}$、纵向场峰值位置和重叠因子；
  \item 对层厚与折射率做小扰动扫描，判断结果是否过于脆弱。
\end{enumerate}

在 COMSOL 路径里，这一阶段通常用二维横截面的 \software{Wave Optics} 本征/模态分析来完成。最小过程是：建立二维层栈几何、指定材料、用适合导模的边界条件和本征求解得到传播常数，再做后处理积分得到重叠与损耗趋势。这里的核心不是软件界面差异，而是输出必须一致：都要产出一个可供下一阶段使用的纵向模式背景，而不仅仅是一张场图。

\section{阶段 B：周期单胞筛选，目标是排除错误结构而不是宣布最终器件}
有了可靠的纵向背景，才进入单胞筛选阶段。这里的任务是回答：在无限周期、开放结构、小信号框架下，哪些 Gamma 点候选模、哪些偏振状态和哪些辐射通道值得继续投入有限器件验证。

这一阶段最好再拆成两个子步骤。先按 \cref{ch:gamma,ch:pwe} 的方法做一次被动无限周期候选模筛选：用 PWEM 建立候选频段、Gamma 点简并关系和主要平面波分量；再把已经缩小后的候选模窗口交给 RCWA 或周期频域模型，去检查开放辐射、上下出光比例和偏振趋势。这样做的好处是，RCWA 不必在过大的参数空间里盲扫，而是围绕已经识别出的候选模做精修。

Lumerical 路径中，若目标是快速扫描单胞参数和辐射分配，优先使用 RCWA；若需要更灵活地看时域响应或近场细节，也可用周期边界下的 FDTD，但成本通常更高。一个最小 RCWA 过程可以概括为：
\begin{enumerate}
  \item 建立单胞几何、上/下介质和层栈；
  \item 设定法向附近的入射波或端口条件，并扫描波长、孔径、刻蚀深度和微扰参数；
  \item 读取反射、透射、衍射效率、局域场增强和上下辐射分配；
  \item 用收敛的傅里叶阶数确认候选模排序与偏振趋势。
\end{enumerate}

COMSOL 路径中，这一步通常用三维单胞频域模型完成：在横向使用相位周期边界，在上下使用端口或开放边界，并通过频率扫或参数扫读取散射量。若要说得再具体一点，就是“用 Wave Optics 的频域单胞模型来做 RCWA 所擅长的那类周期散射判断”，只不过实现框架不是以 RCWA 名义出现，而是以有限元频域周期散射问题出现。

这一阶段的交付物应当是候选模目录，而不是“最终器件最低阈值”的宣称。目录中至少应包含：候选频段、PWEM 给出的 Gamma 点模族和主要傅里叶分量、目标/竞争模的被动损耗排序、偏振趋势、上下出光比例和对工艺扰动的敏感性。

\section{阶段 C：有限器件全波验证，检查无限周期假设到底失效到什么程度}
通过单胞筛选的结构，下一步就要离开无限周期近似，进入有限孔阵模型。这里的任务不再是看单胞散射特征，而是看真实有限阵列中模式是否仍保持正确排序、边缘是否引入新的泄漏通道、远场是否还能维持目标主瓣和偏振状态。

在 Lumerical 路径里，这一步最常见的是 FDTD。最小过程通常是：
\begin{enumerate}
  \item 把候选单胞复制成有限阵列，并加入真实上下层栈；
  \item 在开放方向使用 PML，在必要时利用对称边界降低计算量；
  \item 用脉冲源、偶极源或与目标模匹配的激励源激发结构；
  \item 通过时域衰减、频谱、场分布和近远场变换提取复频率、$Q$、远场和主模排序；
  \item 对阵列尺寸、PML 厚度、网格和仿真时间窗做收敛检查。
\end{enumerate}

在 COMSOL 路径里，这一步常用三维 Wave Optics 的频域或本征频率研究来做。若重点是共振位置和模式形状，可用本征频率研究并在开放边界加 PML；若重点是给定频率下的响应和远场，则更适合频域研究。需要强调的是，COMSOL 的有限元框架在处理复杂材料依赖和后续多物理场联动时很灵活，但三维有限孔阵的计算代价往往不小，因此必须严格控制模型尺寸和收敛验证。

这一阶段的输出至少应包括：有限阵列主模与竞争模的排序、目标模近场与远场、边缘敏感性、以及“无限周期结论在有限阵列中是否仍成立”的判断。

\section{阶段 D：电注入与热建模，开始把阈值增益映射成阈值电流区间}
如果研究目标是器件而不是纯光学结构，那么到这里就必须进入电学与热学。任务不再是找模，而是问：电流怎样进入器件、载流子密度怎样分布、哪里掉压、哪里发热、这些分布怎样改写模增益与模损耗。

在 Lumerical 路径里，常见的分工是用 \software{CHARGE} 解漂移-扩散和泊松方程，用 \software{HEAT} 解热传导。最小过程可以概括为：
\begin{enumerate}
  \item 依据真实外延层和电极定义半导体几何与掺杂；
  \item 在 \software{CHARGE} 中指定接触边界、复合模型和工作电压/电流，求得 $J(\mathbf{r})$ 与 $N(\mathbf{r})$；
  \item 依据 Joule 热和非辐射复合构造热源，送入 \software{HEAT} 得到 $T(\mathbf{r})$；
  \item 把 $N$ 和 $T$ 场转写成折射率修正、增益修正和吸收修正，回写给光学阶段。
\end{enumerate}

在 COMSOL 路径里，对应的是 \software{Semiconductor} 与 \software{Heat Transfer} 接口。几何通常可与光学模型共享，只是在网格、边界和研究类型上要做针对性修改。最小过程是：
\begin{enumerate}
  \item 在 \software{Semiconductor} 接口中设置接触、掺杂、复合与偏置，解出电势和载流子分布；
  \item 在 \software{Heat Transfer} 接口中加入 Joule 热与复合热源，求稳态或瞬态温度场；
  \item 将温度和载流子结果通过自定义材料关系回写到折射率和增益代理中；
  \item 重新运行 \software{Wave Optics}，检查目标模与竞争模的净裕量是否改变。
\end{enumerate}

这里必须再次提醒：除非你真的完成了这个映射链，否则不能把“阈值增益估计值”写成“阈值电流”。

\section{阶段 E：热-电-光迭代，不追求一轮算完，而追求工作窗口收敛}
许多新手会把多物理场耦合理解成“把所有接口同时打开，一次求解结束”。对 PCSEL 来说，这种理解往往既低效也不稳健。更实用的策略通常是分阶段迭代：光学给出模式和重叠，电学给出注入分布，热学给出温度场，再把这些结果回写到光学模型，重复直到目标指标收敛。

工作流真正要追求的，不是某一次迭代里的最佳单点，而是工作窗口是否稳定。一个结构若只有在第一次冷腔迭代中表现良好，而一旦写回温升就发生模排序翻转，它就不能算成功设计。相反，一个方案即使冷腔上没有最漂亮的 $Q$，若热回写后仍保持目标模的净裕量为正，它在器件层面上更有价值。

\section{阶段 F：鲁棒性验证，检查工作窗口是否同时具有工艺窗口}
若阶段 A 到阶段 E 证明的是“名义器件在当前模型下可行”，那么阶段 F 要回答的是“这是不是一个能加工出来、能批量重复出来的器件方案”。也就是说，在名义工作窗口已经建立之后，还应把 \cref{ch:robustness} 中的 tolerance analysis 显式纳入流程。最小动作包括：定义关键误差参数、先做局部灵敏度筛查、对高敏感参数做联合抽样，并用通过率而不是单一样本图来报告结果。

这个阶段之所以不能省略，是因为很多设计在名义值上成功，却在孔半径、刻蚀深度或孔径错位上异常脆弱。若不把阶段 F 写进工作流，所谓“可行设计”往往只是在数值环境里可行，而不是在工艺环境里可行。

\section{阶段 G：实验闭环，不让模型停在“可运行”而要走到“可校准”}
真正能进入研究和器件开发日程表的工作流，还必须再多走一步：把模型接回实验。最小闭环应写成“测量 $\rightarrow$ 参数提取 $\rightarrow$ 分层校准 $\rightarrow$ 预测 $\rightarrow$ 回验”。这里的关键不是把实验当成结果展示，而是把它当成模型可用性的约束。

具体做法通常是：先用截面、几何、电学和热学的直接测量值约束几何与边界条件；再用 I--V、L--I、热阻或谱漂移数据分别校准电学和热学子模型；最后才把这些中间结果回写给光学与增益模型，去预测未参与拟合的阈值漂移、远场变化或偏振窗口。只有当“未参与校准的观测量”也能被解释时，模型才具备真正的外推能力。

这一阶段的更细致写法，会在 \cref{ch:device-appendix} 中系统展开。本章只强调一个纪律：\textbf{不要把所有实验量都同时拿来拟合同一套自由参数。}那样做很容易得到“处处都像、处处都不可信”的黑箱模型。

\section{Lumerical 与 COMSOL 的两条常见实现路径}
前面分阶段讲的是物理链，下面把它压缩成两条常用工程路径，供读者实际启动项目时参考。需要强调，这两条路径都是\textbf{推荐流程}，不是唯一流程。

\subsection*{路径一：以 Lumerical 为主的光学优先链}
这条路径适合“先把单胞和有限阵列光学迅速扫清，再把最有希望的结构接到器件级电热模型”这一工作风格。典型顺序是：
\begin{enumerate}
  \item MODE/FDE：检查纵向模、$n_{\mathrm{eff}}$、$\Gamma_{\mathrm{QW}}$ 与损耗重叠；
  \item RCWA：扫单胞、提取候选模、上下出光和偏振趋势；
  \item FDTD：有限阵列验证、提取远场和边缘敏感性；
  \item CHARGE：求注入分布和串联电阻来源；
  \item HEAT：求温升和热热点；
  \item 脚本化回写：把 $N(\mathbf{r})$ 与 $T(\mathbf{r})$ 回写给光学模型，更新模增益和损耗排序。
\end{enumerate}
这条路径的优点是光学前端筛选效率高，缺点是多物理场耦合通常需要较明确的接口管理和脚本自动化。

\subsection*{路径二：以 COMSOL 为主的多物理场优先链}
这条路径适合“在统一几何和 PDE 框架下做电-热-光之间的参数回写和灵敏度分析”这一工作风格。典型顺序是：
\begin{enumerate}
  \item 二维或简化三维 Wave Optics：建立纵向模背景；
  \item 三维单胞 Wave Optics：做周期单胞散射与偏振筛选；
  \item 三维有限阵列 Wave Optics：做有限器件近场、远场与本征模验证；
  \item Semiconductor：求注入和载流子分布；
  \item Heat Transfer：求温度场与热热点；
  \item 参数回写与分阶段迭代：更新材料参数、重跑光学、检查工作窗口。
\end{enumerate}
这条路径的优点是多物理场接口统一、参数回写自然，缺点是大尺寸三维光学问题的计算成本往往更高，对网格和求解设置更敏感。

\begin{ExampleBox}
若研究目标是快速筛选 50 组单胞设计，再从中挑 2 组做器件级深入验证，通常更适合先走 Lumerical 光学优先链；若研究目标是深入分析一个给定外延与电极方案中，温升如何改变模式竞争，则更适合走 COMSOL 多物理场优先链。这里的“更适合”是工程判断，不是排他性规则。
\end{ExampleBox}

\begin{figure}[htbp]
  \centering
  \begin{tikzpicture}[node distance=0.7cm, every node/.style={draw, rounded corners, align=center, minimum height=0.9cm, text width=3.2cm}]
    \node (a) {阶段 A\\纵向模与层栈检查};
    \node (b) [below=of a] {阶段 B\\周期单胞筛选};
    \node (c) [below=of b] {阶段 C\\有限器件全波验证};
    \node (d) [below=of c] {阶段 D\\电注入与热建模};
    \node (e) [below=of d] {阶段 E\\热-电-光回写与工作窗口};
    \node (l1) [right=2.2cm of a, text width=4.4cm] {Lumerical 风格\\MODE/FDE};
    \node (l2) [right=2.2cm of b, text width=4.4cm] {Lumerical 风格\\RCWA 或周期 FDTD};
    \node (l3) [right=2.2cm of c, text width=4.4cm] {Lumerical 风格\\FDTD};
    \node (l4) [right=2.2cm of d, text width=4.4cm] {Lumerical 风格\\CHARGE + HEAT};
    \node (l5) [right=2.2cm of e, text width=4.4cm] {脚本化接口\\更新 $N,T,\Delta n,\Delta g$};
    \draw[->] (a)--(b);
    \draw[->] (b)--(c);
    \draw[->] (c)--(d);
    \draw[->] (d)--(e);
    \draw[->] (a)--(l1);
    \draw[->] (b)--(l2);
    \draw[->] (c)--(l3);
    \draw[->] (d)--(l4);
    \draw[->] (e)--(l5);
  \end{tikzpicture}
  \caption{PCSEL 仿真工作流示意。图中以 Lumerical 风格链为例标出常见映射；COMSOL 路径遵循同一物理阶段划分，只是在实现上更多使用统一 PDE 框架。}
  \label{fig:workflow-map}
\end{figure}

\section{何时停止迭代：停机判据比“再跑一轮看看”更重要}
完整工作流很容易陷入无止境试错，因此必须在开工前就设定停机判据。最有用的停机条件通常有三类：
\begin{enumerate}
  \item 数值收敛停机：网格、时间窗、PML、傅里叶阶数继续加密时，关键输出变化已低于预设阈值；
  \item 物理一致性停机：功率守恒、参数量纲、层级结论和趋势解释均已自洽；
  \item 目标窗口停机：在目标工作区间内，目标模对竞争模的净裕量、远场和温升都满足要求。
\end{enumerate}
如果这三条没有满足，再漂亮的结果图也不应进入最终结论。

\begin{PitfallBox}
“已经跑出了一个看起来不错的图，所以先写结论再说”是仿真工作流中最常见的失控点。对 PCSEL 来说，真正值钱的不是一张漂亮图，而是一条从参数到账本、从候选模到工作窗口都可追溯的证据链。
\end{PitfallBox}
\section{仿真 - 实验对标协议：如何让数字与测量对话}
仿真的终极价值不在于自洽，而在于能够预测和解释实验。然而，大量 PCSEL 仿真工作在迈向实验对标的第一步就失败了：要么比较的对象不对等，要么忽略了关键的不确定度来源，要么用事后拟合冒充事前预测。本节给出一套严格的仿真 - 实验对标协议，确保每一个"吻合良好"的宣称都经得起推敲。

\subsection{对标前的层级对齐检查}
在将仿真结果与实验数据放在一起比较之前，必须先完成以下五项层级对齐检查：

\begin{ExperimentalAlignmentBox}
\textbf{仿真 - 实验对标的五项前置检查}
\begin{enumerate}
  \item \textbf{几何保真度检查}：仿真中的几何参数（晶格常数$a$、孔径$r/a$、刻蚀深度$h_{\text{etch}}$、外延层厚）是否与 SEM/TEM 实测一致？若使用名义设计值而非实测值，必须说明理由并评估敏感度。
  
  \item \textbf{材料模型真实性检查}：折射率和增益谱是来自实测椭偏/PL 数据，还是仅靠文献经验公式？对有源区，是否包含了载流子浓度依赖的增益压缩和折射率变化？
  
  \item \textbf{边界条件等价性检查}：仿真的是无限周期单胞、有限尺寸阵列还是完整器件？实验测得的是面发射主瓣、边缘发射还是积分球总功率？两者是否在相同的物理量层面进行比较？
  
  \item \textbf{工作点对齐检查}：仿真的是阈值处、阈值以上固定倍数还是最大功率点？实验数据是在 CW 还是脉冲条件下测得？若实验未明确热沉温度和驱动条件，仿真不应直接声称"绝对阈值电流吻合"。
  
  \item \textbf{不确定度量化检查}：仿真侧是否提供了收敛误差估计？实验侧是否给出了仪器精度、样品间涨落和重复性误差？若任何一方缺失不确定度信息，比较时应降级为"定性趋势对比"而非"定量验证"。
\end{enumerate}
\end{ExperimentalAlignmentBox}

\begin{WhatCouldGoWrong}
最常见的对标失效案例是：仿真给出的是冷腔$Q$值换算的"阈值增益"，实验给出的是 LIV 曲线上的转折点对应的"注入电流"。这两个量虽然相关，但中间隔着载流子输运效率、内量子效率、串联电阻和非辐射复合等多个黑箱。直接将二者画在同一张图上并声称"验证了阈值模型"，属于典型的范畴错误。
\end{WhatCouldGoWrong}

\subsection{三级对标体系：从定性到定量}
根据可用信息的完整程度和对标目标的差异，建议采用三级对标体系：

\paragraph{Level 1: 定性趋势对标（最低门槛）}
这一级只要求仿真正确预测参数变化的方向，不要求幅度准确。典型问题包括：
\begin{itemize}
  \item 孔径增大时，阈值是先降后升还是单调变化？
  \item 刻蚀加深时，远场主瓣变窄还是变宽？
  \item 温度升高时，激射波长红移速率是多少 nm/K？
\end{itemize}
通过标准：仿真预测的趋势符号与实验一致，且不存在反例。

\begin{ReportingStandardBox}
\textbf{Level 1 对标的规范表述示例}
"如图 X 所示，仿真预测随着占空比从 0.3 增加到 0.6，A 模的辐射效率单调上升，这与 PL 淬灭校正后的 EL 强度趋势一致（Pearson 相关系数$r=0.87$）。需强调这是定性趋势验证，不涉及绝对效率标定。"
\end{ReportingStandardBox}

\paragraph{Level 2: 半定量范围对标（推荐门槛）}
这一级要求仿真给出的数值落在实验测量的不确定度范围内，或至少在同一数量级。典型问题包括：
\begin{itemize}
  \item 仿真阈值电流密度$J_{\text{th,sim}}$是否处于实验值的$\pm 30\%$以内？
  \item 远场发散角的仿真值与 CCD 拟合值偏差是否小于$5^\circ$？
  \item 模式波长差的仿真预测与光谱仪分辨率相比是否显著？
\end{itemize}
通过标准：关键量的相对偏差$\delta = |X_{\text{sim}}-X_{\text{exp}}|/X_{\text{exp}} < 30\%$，且所有偏差均可由已知不确定度源解释。

\begin{ExperimentalAlignmentBox}
\textbf{Level 2 对标的不确定度预算模板}
若仿真值为$X_{\text{sim}} \pm \sigma_{\text{sim}}$，实验值为$X_{\text{exp}} \pm \sigma_{\text{exp}}$，则兼容性判据为：
\[
|X_{\text{sim}} - X_{\text{exp}}| \leq k \sqrt{\sigma_{\text{sim}}^2 + \sigma_{\text{exp}}^2},
\]
其中$k$是覆盖因子（通常取$k=2$对应 95\%置信区间）。若上述不等式不成立，则需要检查是否存在系统性偏差（如材料参数失配或未建模物理机制）。
\end{ExperimentalAlignmentBox}

\paragraph{Level 3: 定量预测对标（严格门槛）}
这一级要求在没有事后拟合的前提下，仿真预测与实验数据在误差棒内完全吻合。这通常需要：
\begin{itemize}
  \item 所有输入参数均来自独立表征（SEM、椭偏、Hall 效应、DLTS 等）
  \item 仿真模型包含所有已知的一阶和二阶物理机制
  \item 进行了盲测验证：用一部分实验数据校准模型，然后用另一部分数据进行预测检验
\end{itemize}
通过标准： reduced chi-square $\chi_\nu^2 = \frac{1}{N-p}\sum_i \frac{(y_{\text{sim},i}-y_{\text{exp},i})^2}{\sigma_i^2}$接近 1（理想范围 0.8--1.2），其中$N$是数据点数，$p$是拟合参数个数。

\begin{WhatCouldGoWrong}
Level 3 对标的一个常见陷阱是"过度调参"：通过调整多个自由参数（如有源区掺杂浓度、缺陷态密度、接触电阻等）使得仿真曲线与实验点完美重合，但这些参数的物理合理性未经独立验证。这种做法得到的"吻合"只是数学拟合，不具备预测能力。正确的做法是报告每个可调参数的物理约束范围和敏感性分析。
\end{WhatCouldGoWrong}

\subsection{关键物理量的对标技术规范}
针对 PCSEL 特有的几类关键测量，以下是具体的对标技术要点：

\subsubsection{LIV 曲线的对标}
LIV（Light-Current-Voltage）特性是激光器最基本的宏观指标。对标时要注意：

\begin{itemize}
  \item \textbf{阈值提取方法一致性}：实验上常用拐点法、线性拟合法或二次微分法提取$I_{\text{th}}$，仿真上也应采用相同算法处理$L-I$曲线，而不是直接用冷腔$Q$值推算。
  
  \item \textbf{热滚降（thermal rollover）建模}：若实验观察到高温或大电流下的功率饱和/下降，仿真必须包含自加热效应和温度依赖的效率下降机制（如 Auger 复合、载流子泄漏）。
  
  \item \textbf{串联电阻校准}：$V-I$曲线的斜率应与实验 extracted 的串联电阻$R_s$一致。若偏差超过$20\%$，需要检查接触模型和掺杂分布。
\end{itemize}

\subsubsection{光谱特性的对标}
激射波长和边模抑制比（SMSR）的对标要点：

\begin{itemize}
  \item \textbf{波长温度漂移}：实验测得的$\dd\lambda/\dd T$应与仿真中材料色散 + 热光效应的综合预测一致。典型 InGaAsP/InP体系的$\dd\lambda/\dd T \approx 0.05\text{--}0.1\,$nm/K。
  
  \item \textbf{SMSR的定义统一}：实验 SMSR 通常定义为主模与最强边模的功率比（dB），仿真也应在相同带宽和积分时间内计算，避免因频谱分辨率不同导致的人为高估。
  
  \item \textbf{动态展宽效应}：若实验光谱仪分辨率有限（如$>0.1\,$nm），仿真应考虑线宽展宽机制（相位噪声、模式跳动），否则可能出现"仿真线宽远窄于实验"的假性不符。
\end{itemize}

\subsubsection{远场图案的对标}
远场辐射图案的对标最具挑战性，因为其对几何误差极为敏感：

\begin{itemize}
  \item \textbf{角度坐标系统一}：确认实验使用的是立体角$(\theta,\phi)$还是直角坐标$(x,y)$探测器平面，并进行适当的 Jacobian 变换。
  
  \item \textbf{背景扣除}：实验远场常包含杂散光和基底辐射，对标前应减去暗背景和自发辐射背景。
  
  \item \textbf{有限尺寸效应}：若仿真用的是无限周期近似，而实验器件只有几十个周期的阵列，衍射旁瓣的位置和强度会有显著差异。此时应改用有限尺寸全波仿真或直接使用 Fraunhofer 衍射公式修正。
\end{itemize}

\begin{CrossValidationBox}
\textbf{远场对标的快速诊断流程}
\begin{enumerate}
  \item 比较主瓣 FWHM：若仿真值明显小于实验值，可能是阵列尺寸过大或孔径均匀性假设过理想
  \item 比较旁瓣位置：若旁瓣角度系统性偏移，可能是晶格常数或波长输入有误
  \item 比较偏振消光比：若实验 PER 远低于仿真值，可能存在应力双折射或加工不对称
  \item 旋转平均 vs 方位分辨：确认实验是否做了角度平均，避免将各向异性仿真与各向同性实验直接比较
\end{enumerate}
\end{CrossValidationBox}

\subsection{当仿真与实验不符时的诊断树}
即使完成了上述所有检查，仿真与实验出现分歧仍是常态。这时应避免两种极端反应：一是盲目怀疑仿真，二是轻易否定实验。建议采用如下系统化诊断流程：

\begin{enumerate}
  \item \textbf{复查输入参数}：逐一核对所有几何和材料参数，特别注意单位制转换（nm/$\mu$m/eV/meV）、温度条件和批次间工艺变异。
  
  \item \textbf{简化测试}：构建一个极度简化的基准案例（如单层平板波导、理想 PEC 边界），验证求解器本身是否正确配置。
  
  \item \textbf{网格/基函数收敛再确认}：有时初始收敛测试不够充分，特别是在引入新物理场（如热耦合）后需要重新验证。
  
  \item \textbf{敏感性分析}：对可疑参数做扫描，观察是否存在某个参数的微小变化就能解释仿真 - 实验差距。这可能揭示未被充分认识的关键机制。
  
  \item \textbf{独立工具复核}：用另一种数值方法（如 FDTD vs FEM）或另一款软件重算同一结构，排除单一工具的 bugs 或实现细节影响。
  
  \item \textbf{请求原始数据}：若实验数据来自合作组，索要原始测量文件和校准记录，有时数据处理过程中的平滑、滤波或归一化会引入人为 artifacts。
\end{enumerate}

\begin{PracticalTipBox}
建立一个"失败日志"文档，详细记录每一次仿真 - 实验不符的案例、排查过程和最终原因。长期积累下来，你会发现某些类型的偏差反复出现（如始终低估阈值电流 20\%），这往往指向某个被忽视的系统性误差源（如接触电阻模型过于理想）。
\end{PracticalTipBox}

\subsection{对标结果的报告伦理}
最后，必须强调对标报告的学术诚信底线：

\begin{EthicsReminderBox}
\textbf{仿真 - 实验对标的红线准则}
\begin{itemize}
  \item \textbf{禁止选择性呈现}：不得只展示吻合良好的数据点而隐藏明显偏离的点。若某些条件下的对标失败，应如实报告并讨论可能原因。
  
  \item \textbf{禁止事后调参冒充预测}：若使用了实验数据来确定模型参数，必须声明这是"拟合"而非"预测"，并用独立的测试集进行验证。
  
  \item \textbf{禁止模糊不确定度}：不得用光滑曲线连接稀疏的实验点制造"完美跟踪"的视觉错觉。应清晰标示误差棒和数据点密度。
  
  \item \textbf{禁止跨层级比较}：不得将冷腔仿真结果直接与器件级 LIV 曲线放在同一纵轴上比较，除非明确说明了中间的转换假设。
\end{itemize}
\end{EthicsReminderBox}
\section*{本章小结}
一套实用的 PCSEL 仿真工作流必须以结论分层和参数账本为起点，然后依次经历纵向模检查、周期单胞筛选、有限器件验证、电学与热学建模，以及热-电-光回写。Lumerical 与 COMSOL 都能承载这条链，但常见工程分工不同：前者通常更擅长高效光学筛选，后者通常更擅长统一多物理场耦合。无论选哪条路线，真正应追求的都不是“某一轮最优图片”，而是一个通过收敛和物理一致性检查的工作窗口。若把本章和后面两章一起看，本章给出的是正向流程骨架，\cref{ch:worked-example} 给出的是一次完整走法，\cref{ch:failure-modes} 给出的则是反向诊断顺序。

\section*{练习题}
\begin{enumerate}
  \item \basicq 为一个你感兴趣的 PCSEL 方案写出 L1 到 L4 的结论层级，并说明每一级允许说什么、不允许说什么。

  \item \basicq 分别给出阶段 A 到阶段 E 的输入与输出变量，特别指出哪些量需要在不同求解器之间传递。

  \item \midq 比较以 Lumerical 为主和以 COMSOL 为主的两条工作流，说明它们各自最可能在哪一步最耗时、最容易出错。

  \item \hardq 为你的项目设计三条停机判据，并解释为什么这些判据能防止你过早下结论。
\end{enumerate}

\section*{建议继续阅读}
\begin{itemize}
  \item 下一章将把本章工作流落实到一个完整示例，给出参数、建模顺序、求解设置模板和可信度检查。阅读时建议对照 \cref{tab:workflow-example-failure}，逐步核对每一阶段到底在例题中对应哪一步。 这条阅读的重点是把本章的输出顺着主线接到后续模块，而不是停成孤立知识点。

  \item 结合 \cref{ch:failure-modes} 阅读，可以把本章工作流直接转化为项目级排错流程；最有效的用法，是先按本章正向走，再按失败模式章反向复盘。 这条阅读的重点是把本章的输出顺着主线接到后续模块，而不是停成孤立知识点。
\end{itemize}

